\section{Prompts and tool schemas, verbatim}
\label{app:prompts}

All prompts are reproduced exactly as sent, in the original Chinese. We do not translate
them: the protocol comparison is between two prompts, and a translation is a third prompt.
An English gloss follows each block for readers who do not read Chinese, but the gloss was
never sent to any model.

\subsection{ReAct}

Two strengths. \texttt{optional} permits a direct answer; \texttt{mandatory} requires an
\texttt{Action} first. The main table uses \texttt{optional} on both protocols so that
neither side is coerced.

{\footnotesize
\begin{verbatim}
REACT_OPTIONAL
--------------
你是一个 Python 编程助手，可以按以下格式工作：

Thought: 你的思考
Action: run_tests
Action Input: ```python
<完整代码>
```
Observation: （由系统填写测试结果）

Thought: 我已确认代码正确
Final Answer: ```python
<最终代码>
```

你可以用 Action 提交代码验证，也可以直接给 Final Answer。
\end{verbatim}
}

\noindent\emph{Gloss:} ``You are a Python coding assistant and may work in the following
format: Thought / Action: run\_tests / Action Input: (fenced code) / Observation (filled in
by the system) / \ldots{} / Final Answer: (fenced code). You may submit code with Action for
verification, or give a Final Answer directly.''

The \texttt{mandatory} variant is identical except that the closing line becomes
\textbf{必须先用 Action 提交代码验证，再给 Final Answer。}
(``You must first submit code for verification with Action, then give a Final Answer''), and
the loop is marked as repeatable.

\subsection{Function calling}

{\footnotesize
\begin{verbatim}
FC_OPTIONAL
-----------
你是一个 Python 编程助手。请根据用户的描述实现所要求的函数。

你可以调用 run_tests 工具把代码交给测试运行，工具会返回通过情况与报错，
你可以据此修改代码再次提交。

FC_MANDATORY
------------
你是一个 Python 编程助手。请根据用户的描述实现所要求的函数。

**你必须先调用 run_tests 工具验证你的代码**，不要直接给出答案。
工具会返回测试通过情况与报错；若有失败，请修改代码并再次调用 run_tests。
\end{verbatim}
}

\noindent\emph{Gloss:} optional --- ``You may call the run\_tests tool to submit code to the
test runner; it returns pass/fail and errors, and you may revise and resubmit.''
mandatory --- ``\textbf{You must first call the run\_tests tool to verify your code}; do not
answer directly.''

\textbf{Strength is matched across protocols in every comparison we report.} The training
comparison in \S\ref{sec:rl} uses \texttt{mandatory} on both sides, so ``FC executed zero tool calls''
cannot be attributed to a weaker instruction.

\subsection{Tool schemas: terse, rich, strict}

\texttt{terse} is the schema a reader arrives at by following the API documentation. It is
the schema under which Llama's 23\% wrong-tool rate appears.

{\footnotesize
\begin{verbatim}
terse
-----
{"type": "function", "function": {
  "name": "run_tests",
  "description": "把你写的 Python 代码交给测试运行，返回通过情况与报错。",
  "parameters": {"type": "object",
                 "properties": {"code": {"type": "string"}},
                 "required": ["code"]}}}
\end{verbatim}
}

\texttt{rich} is the \S\ref{sec:llama} control. Its purpose is to test whether the wrong-tool rate is
simply an under-specified schema, so it is written to be maximally explicit: the description
states that the tool is \emph{not} the function the user asked for, and the parameter
description carries a worked example.

{\footnotesize
\begin{verbatim}
rich
----
{"type": "function", "function": {
  "name": "run_tests",
  "description": "把你写的 Python 代码提交给测试运行器执行，返回测试通过情况与报错。"
                 "注意：这个工具不是你要实现的那个函数，不要把题目函数的参数传进来；"
                 "唯一的参数 code 是你写的完整源代码文本。",
  "parameters": {"type": "object",
    "properties": {"code": {
      "type": "string",
      "description": "完整的 Python 源代码文本，包含所需的 import 和完整的"
                     "函数定义。例如：\"def add(a, b):\\n    return a + b\""}},
    "required": ["code"]}}}
\end{verbatim}
}

\noindent\emph{Gloss of the added text:} ``Note: this tool is not the function you are being
asked to implement. Do not pass the task function's parameters into it. Its only parameter,
\texttt{code}, is the complete source text you wrote.''

\textbf{Under this explicit warning 22/100 items still pass the task function's parameters},
and they are the same parameter names that failed under \texttt{terse} (\S\ref{sec:llama}, $p$=1.000).

\texttt{strict} is \texttt{terse} plus \verb|"strict": true| on the function object, with
\texttt{VLLM\_ENFORCE\_STRICT\_TOOL\_CALLING=true} (the default). This is the single change
that takes the wrong-tool rate to zero.

\subsection{The Thought-scaffold control (\S\ref{sec:llama})}

To test whether ReAct's advantage is merely an extra reasoning step, the FC system prompt was
given a Thought requirement and nothing else changed:

\begin{quote}
``First write out your reasoning in a Thought: what the problem asks for, what the edge cases
are, what algorithm you intend to use. Once you have thought it through, call the
\texttt{run\_tests} tool\ldots{}''
\end{quote}

The wrong-tool rate rose from 23 to 59 ($p<0.001$). We report this as a single wording, not
as a claim about reasoning scaffolds in general.

\subsection{The role-disambiguation few-shot (\S\ref{sec:rl})}

Used at training scale (1.5B) after the wrong-tool failure was found to be semantic rather
than syntactic. It shows the wrong call and the right call side by side.

{\footnotesize
\begin{verbatim}
你是一个 Python 编程助手。请根据用户的描述实现所要求的函数。

**你必须先调用 run_tests 工具验证你的代码**，不要直接给出答案。

注意一个常见错误：**用户描述里要你实现的那个函数，不是可调用的工具。**
你唯一可以调用的工具是 `run_tests`，它只有一个参数 `code`，内容是你写的完整源代码。

错误示范（用户要求实现 join_integers）：
{"name": "join_integers", "arguments": {"int_array": [1, 2, 3]}}
—— 这是在调用你自己该编写的函数，run_tests 收不到任何代码。

正确示范：
{"name": "run_tests", "arguments": {"code": "def join_integers(int_array):\n
   return ','.join(map(str, int_array))"}}
—— 函数名恒为 run_tests，你写的整段代码作为字符串放进 code。

工具会返回测试通过情况与报错；若有失败，请修改代码并再次调用 run_tests。
\end{verbatim}
}

\textbf{It made things worse}: recoverable calls rose 52 $\to$ 64 while calls naming
\texttt{run\_tests} stayed at \textbf{0} and final pass fell 15 $\to$ 13. The model became
more fluent at emitting the wrong call.

\subsection{The intent criterion}

Every emitted-call count in this paper comes from one function, applied identically to every
arm. There is no per-arm tuning.

{\footnotesize
\begin{verbatim}
TIGHT = re.compile(
    r'"name"\s*:\s*"run_tests".{0,200}?"arguments"\s*:\s*\{'
    r'.{0,80}?"code"\s*:\s*"(.{0,4000}?)"\s*\}', re.S)

REAL  = re.compile(r'\\n|def |class |return |import |lambda ')

def is_tight_call(raw):
    m = TIGHT.search(raw or "")
    return bool(m and REAL.search(m.group(1)))
\end{verbatim}
}

\texttt{TIGHT} requires the name, an \texttt{arguments} object, and a \texttt{code} key whose
value is a \emph{string literal}; \texttt{REAL} then requires that literal to contain actual
Python. Together they exclude schema echo (\texttt{parameters}/\texttt{description} present,
\texttt{arguments} absent), identifiers passed instead of strings
(\texttt{"code": test\_code}), and placeholders (\texttt{"<your code here>"}).

The \texttt{strong} tier drops the payload requirement (name present, or an XML tool tag);
\texttt{weak} counts JSON structure with no \texttt{run\_tests} name and is disjoint from
both. Tiers are nested, not mutually exclusive, and the paper states which tier each number
uses. Human validation of this criterion is in \S3.2; two independent blind annotators give
inter-annotator $\kappa$ = 0.967 on items where both read identical bytes.
